\documentclass[10pt,a4paper]{article}

\usepackage[top=0.48in,bottom=0.48in,left=0.62in,right=0.62in]{geometry}
\usepackage[hidelinks]{hyperref}
\usepackage{enumitem}
\usepackage{titlesec}
\usepackage{array}
\input{glyphtounicode}
\pdfgentounicode=1
\ifdefined\pdfinterwordspaceon
  \pdfinterwordspaceon
\fi

\pagestyle{empty}
\setlength{\parindent}{0pt}
\setlength{\parskip}{0pt}
\linespread{1.00}

\setlist[itemize]{
  label=-,
  leftmargin=1.15em,
  itemsep=0pt,
  topsep=1pt,
  parsep=0pt,
  partopsep=0pt
}

\titleformat{\section}
  {\large\bfseries}
  {}
  {0pt}
  {}
  [\vspace{-4pt}\rule{\linewidth}{0.45pt}]

\titlespacing{\section}{0pt}{5pt}{2pt}

\newcommand{\entry}[2]{
  \vspace{2pt}
  \begin{tabular*}{\textwidth}{@{}l@{\extracolsep{\fill}}r@{}}
    \textbf{#1} & #2 \\
  \end{tabular*}
}

\newcommand{\subentry}[1]{\textit{#1}\vspace{-1pt}}

\newcommand{\skillrow}[2]{
  \begin{tabular*}{\textwidth}{@{}p{0.23\textwidth}@{\extracolsep{\fill}}p{0.73\textwidth}@{}}
    \textbf{#1} & #2 \\
  \end{tabular*}
}

\begin{document}

\begin{center}
{\LARGE \textbf{Gioia Zheng}} \\[3pt]
Rome, Italy \quad | \quad
\href{mailto:zheng.2040439@studenti.uniroma1.it}{zheng.2040439@studenti.uniroma1.it} \quad | \quad
\href{https://www.linkedin.com/in/gioia-zheng-9233a0303}{linkedin.com/in/gioia-zheng-9233a0303} \\[2pt]
\href{https://github.com/GioiaZheng}{github.com/GioiaZheng} \quad | \quad
\href{https://gioiazheng.github.io}{gioiazheng.github.io} \quad | \quad
\href{https://orcid.org/0009-0000-8212-4506}{ORCID: 0009-0000-8212-4506}
\end{center}

\vspace{-5pt}

\section*{Research Profile}
Applied Computer Science and AI undergraduate at Sapienza University of Rome working on safety-relevant evaluation of RAG and LLM systems. I build reproducible pipelines for paired comparisons, grounding checks, trace-level diagnostics, and failure taxonomies that expose regressions hidden by aggregate metrics.

\section*{Education}
\entry{Sapienza University of Rome}{Sep 2022 -- Present}
B.Sc. in Applied Computer Science and Artificial Intelligence

\section*{Research Experience}
\entry{NLP Algorithm Engineer Intern, Microsoft (China)}{Feb 2026 -- May 2026}
\subentry{Natural Language Computing, remote}
\begin{itemize}
\item Designed and implemented a multi-stage retrieval and generation evaluation pipeline over MS MARCO-scale data: BM25 retrieval, FAISS dense retrieval, transformer reranking, and seq2seq generation.
\item Built paired evaluation routines with confidence intervals to measure how stage-level retrieval changes propagate to answer quality and grounding behavior.
\item Defined schema-versioned experiment manifests and validation checks so results remain traceable to code, data, configuration, and environment.
\end{itemize}

\section*{Selected Research Projects}
\entry{MS MARCO Generative QA System}{\href{https://github.com/GioiaZheng/msmarco-genqa}{github.com/GioiaZheng/msmarco-genqa}}
\subentry{Python, PyTorch, BM25, FAISS, cross-encoder reranking, T5, grounding evaluation}
\begin{itemize}
\item Study when retrieval and reranking gains translate into faithful, evidence-supported generation rather than only higher aggregate surface metrics.
\item Built a reproducible retrieve $\rightarrow$ rerank $\rightarrow$ generate pipeline on MS MARCO passage data with config-driven experiments and schema-versioned manifests.
\item Reported Token-F1 improvement from 0.197 to 0.368 on 6,980 paired queries, with paired-bootstrap confidence interval [+0.163, +0.178].
\item Developed a failure taxonomy separating retrieval, reranking, evidence-use, and generation errors to expose regressions hidden by aggregate scores.
\end{itemize}

\entry{RAG Observatory}{\href{https://github.com/GioiaZheng/rag-observatory}{github.com/GioiaZheng/rag-observatory}}
\subentry{Python, trace schemas, failure taxonomy, evidence attribution, diagnostic reports}
\begin{itemize}
\item Building a trace-based observability toolkit for RAG runs covering queries, retrieved candidates, context, answers, evidence, evaluation signals, and failure labels.
\item Implemented schema validation, failure-label checks, diagnostic reports, and comparison workflows for inspecting support, abstention behavior, and metric-hidden regressions.
\end{itemize}

\section*{Additional Experience}
\entry{Web Development Intern, Idearia}{Mar 2026 -- Jun 2026}
Maintain production systems and automate operational workflows in a client-facing environment.

\entry{Young Representative, UPI Emilia-Romagna}{Sep 2022 -- Oct 2023}
Conducted data-driven research on school dropout and contributed to policy-oriented reports.

\section*{Skills}
\skillrow{Programming}{Python, SQL}
\skillrow{ML and Retrieval}{Information retrieval, RAG and LLM evaluation, PyTorch, FAISS, BM25, Sentence Transformers, cross-encoder reranking, T5}
\skillrow{Research Methods}{Paired bootstrap, confidence intervals, controlled ablations, grounding evaluation, failure taxonomy, trace-level diagnostics}
\skillrow{Research Tooling}{Linux, Docker, Git, GitHub Actions, experiment manifests, reproducible pipelines}
\skillrow{Languages}{Chinese (Native), Italian (Bilingual), English (Professional), German (Beginner)}

\end{document}
